\documentclass[10pt,a4paper]{amsart}
\usepackage[margin=19mm]{geometry}
\usepackage{amsmath,amssymb,mathtools}
\usepackage[scr=rsfs,cal=euler]{mathalfa}
\usepackage{xfrac}
\usepackage[unicode,hidelinks,bookmarks=false]{hyperref}
\newcommand{\ie}{i.e.\ }
\newcommand{\cf}{cf.\ }
\newcommand{\resp}{respectively\ }
\newcommand{\Subsection}{\subsection}
\providecommand{\nobreakdash}{\nobreak\hbox{-}}
\setlength{\emergencystretch}{2em}
\multlinegap0pt
\usepackage{mdframed}
\usepackage{mdframed}
\usepackage{mdframed}
\usepackage{mdframed}
\usepackage{amssymb}
\usepackage{enumerate}
\usepackage{mathtools}
\usepackage{mdframed}
\usepackage[shortlabels]{enumitem}
\newtheorem{proposition}{Proposition}
\DeclareMathOperator{\Hom}{Hom}
\newcommand{\Spec}{{\rm Spec }}
\DeclareMathOperator{\Tr}{Tr}
\DeclareMathOperator{\ide}{Id}
\title{Transfer and Norm for Finite Group Schemes}
\author{Kostas Karagiannis, Peter Symonds}
\date{Source: arXiv:2603.28501v2 (CC BY 4.0)}
\begin{document}
\maketitle

\noindent\emph{Source and licence.} Transfer and Norm for Finite Group Schemes. Version arXiv:2603.28501v2 (CC BY 4.0), distributed by arXiv under the Creative Commons Attribution 4.0 International licence, http://creativecommons.org/licenses/by/4.0/, as recorded in the arXiv licence metadata for this exact version and shown on the abstract page https://arxiv.org/abs/2603.28501v2. Full licence notice: \texttt{licenses/arxiv-2603.28501-CC-BY-4.0.txt}.\par\smallskip

\noindent\textit{Editorial scope.} Counted unit: the article's proposition pr:properties-tr (source0.tex lines 287--309). The article's complete original proof of it is delivered with it (source0.tex lines 311--319, one \texttt{proof} environment of 9 lines). No mathematics is written, repaired or re-derived: the statement and the proof are the article's own lines, in the article's own notation, and no retained line is substituted or re-flowed.  No further source-local context is needed: every cross-reference of every retained line resolves inside the two delivered spans.  Nothing was omitted from any retained span, so the package carries no omission record.  No retained line contains a citation, so the wrapper carries no bibliography and no bibliography entry.  No retained line contains a figure environment, an image-inclusion command or drawing code, so no image is delivered and the package declares no asset dependency.  Editorial formatting only, in addition: an AMS \texttt{amsart} preamble that restates the article's own declarations and macros used by the retained lines, namely the environments \texttt{proposition} on separate counters, together with the standard packages those lines need.  The retained environments print in this document as Proposition 1 in order of appearance, whereas the article prints the same items with the numbering of its own full document.  The complete native source of the article as distributed in the arXiv e-print for this exact version is bundled unchanged as \texttt{sources/197449/source0.tex}.
\begin{proposition}\label{pr:properties-tr}
Let $L,L',M',M,N',N$ be $G$-modules, and let $K,H$ be closed subgroup schemes of $G$ with $K\leq H$.

\begin{enumerate}
\item 
\begin{enumerate}
	\item \label{tr1a}
	If $d \in \Hom_G(L,L')$ and $ s \in \Hom_H(\omega^{-1}_{H,G}, L)$ then $\Tr^G_H (d \circ s) = d \circ \Tr^G_H(s)$.
	\item \label{tr1b}
If $f \in \Hom_H (M,  \omega_{H,G} \otimes N), \; g \in \Hom_G (M',M)$ and $h \in \Hom_G (N,N')$, then 
\[
h \circ \Tr_H^G(f)\circ g = \Tr_H^G( ( \ide_{\omega_{H,G}} \otimes h) \circ f \circ g).
\]
\end{enumerate}

\item \label{tr2} $\tilde{\Tr}^G_K = \tilde{\Tr}^G_H \circ \tilde{\Tr}^H_K \! :  {}^K({}_K \delta  \otimes L) \rightarrow {}^G({}_G \delta  \otimes L)$ and 
$\Tr^G_K = \Tr^G_H \circ (\ide _{\omega_{H,G}} \otimes \Tr^H_K) \! :  {}^K( \omega_{K,G}  \otimes L) \rightarrow {}^GL$. 

\item \label{tr3} If $\ell /k$  is a field extension, and $G_\ell=G\times_{\Spec (k)} \Spec (\ell),\;H_\ell=H\times_{\Spec (k)} \Spec (\ell)$, then $  \Tr_{H_\ell}^{G_\ell} = \Tr_{H}^G \otimes \ell \! : 
{}^H(\omega_{H_\ell,G_\ell} \otimes L) \rightarrow {}^G(\ell \otimes L)$.

\end{enumerate}
\end{proposition}

\begin{proof}
	
	Part \ref{tr1a} follows from the definitions, since the definition of $\Tr$ only uses the first argument of $\Hom$. Part \ref{tr1b} then follows by setting $L= \Hom(M,N)$ and $L'=\Hom(M',\omega_{H,G}N')$.
	
	Part \ref{tr2} is the observation that the composite
	\[ {}_G \delta \xhookrightarrow{\tilde{t}_{H,G}} kG \otimes _H {}_H \delta \xhookrightarrow{\ide_{kG} \otimes \tilde{t}_{H,G}} kG \otimes_{kH}kH \otimes _{kK} {}_K \delta \cong kG \otimes_K {}_K \delta \]
	is just $\tilde{t}_{K,G}$.
	Finally,  part \ref{tr3} is just the observation that everything we have done commutes with base change.
\end{proof}

\end{document}
